\documentclass[11pt,a4paper]{article}
\usepackage{fontspec,amsmath,amssymb,amsthm,microtype}
\setmainfont{lmroman10-regular.otf}[BoldFont=lmroman10-bold.otf,ItalicFont=lmroman10-italic.otf,BoldItalicFont=lmroman10-bolditalic.otf]
\usepackage{xcolor}
\usepackage[most]{tcolorbox}
\usepackage{mathrsfs}
\providecommand{\lt}{<}
\providecommand{\gt}{>}
\setmonofont{DejaVuSansMono.ttf}[BoldFont=DejaVuSansMono-Bold.ttf,ItalicFont=DejaVuSansMono-Oblique.ttf,BoldItalicFont=DejaVuSansMono-BoldOblique.ttf,Scale=0.86]
\newfontfamily\leanfallback{STIXTwoMath-Regular.otf}[Scale=0.86]
\definecolor{leanrule}{gray}{0.62}
\definecolor{leanback}{gray}{0.962}
\definecolor{leantitle}{gray}{0.22}
\definecolor{leanmuted}{gray}{0.38}
\newtcolorbox{leanbox}[1]{enhanced,breakable,colback=leanback,colframe=leanrule,
  boxrule=0.4pt,arc=0.6mm,left=2.4mm,right=2.4mm,top=1.8mm,bottom=1.8mm,
  before skip=0.8em,after skip=1em,fonttitle=\small\bfseries,coltitle=leantitle,
  colbacktitle=leanback,titlerule=0.3pt,title={#1}}
\newlength{\leancharwidth}
\newenvironment{leancode}{\par\smallskip\begingroup\ttfamily\footnotesize\raggedright
  \settowidth{\leancharwidth}{M}\setlength{\parskip}{0pt}\setlength{\parindent}{0pt}
  \frenchspacing}{\par\endgroup\smallskip}
\newcommand{\leanline}[2]{\hangindent=\dimexpr#1\leancharwidth+4\leancharwidth\relax
  \hangafter=1\noindent\hspace*{#1\leancharwidth}#2\par}
\newcommand{\leansource}[1]{\par{\footnotesize\color{leanmuted}#1}\par}
\newcommand{\leanlabel}[1]{\par\smallskip\noindent{\small\bfseries #1}\par}
\usepackage[margin=25mm]{geometry}
\usepackage[unicode,hidelinks]{hyperref}
\hypersetup{pdfauthor={LangevinLean Project}}
\setlength{\emergencystretch}{3em}
\setlength{\parskip}{0.4em}
\setlength{\parindent}{0pt}
\allowdisplaybreaks[2]
\raggedbottom
\title{\texorpdfstring{Accumulation of bounded-test errors and invariance of a limiting law}{Accumulation of bounded-test errors and invariance of a limiting law}}
\author{LangevinLean Project}
\date{}
\begin{document}\maketitle
Draft mathematical exposition; these notes may contain mistakes.


Corrected invariance can be used entirely through bounded tests. Composing a test with more deterministic proposal steps preserves its bound. Each telescoping increment is therefore controlled under the same fixed measure, giving a linear accumulation of rejection mass. Pointwise convergence and vanishing accumulated error then identify the limiting law exactly by dual characteristic functions.



Telescoping bounded-test errors produces linear accumulation over a fixed horizon. Once the numerical maps converge, characteristic-function uniqueness identifies the invariant law of the limiting Hamiltonian flow.


\section{\texorpdfstring{Telescope tests composed with deterministic iterates}{Telescope tests composed with deterministic iterates}}


Let \(A_k=\int F(V^kz)\,d\pi\). Apply the one-step weak estimate to \(G=F\circ V^k\). Measurability of iterates proves measurability of \(G\), and the same bound \(\|G\|\le C\) holds. Thus \(\|A_{k+1}-A_k\|\le2C\int(1-a)\,d\pi\). The triangle inequality and induction give \[\|A_n-A_0\|\le2Cn\int(1-a)\,d\pi.\] The argument explicitly rewrites the successor iterate as \(V^k\circ V\), applies the one-step result, and uses that the successor step count is one more than the preceding count. Every increment is integrated against the fixed measure \(\pi\); no law of an unadjusted intermediate chain is assumed invariant.



Write the difference between the \(n\)-step test integral and its invariant value as a sum of \(n\) one-step differences. Composition with deterministic iterates preserves the test’s uniform bound. Each term is therefore bounded by the same local rejection estimate, giving linear rather than exponential accumulation.

\begin{leanbox}{Lean: \texttt{correctedMeasure\_\allowbreak{}iterate\_\allowbreak{}weak\_\allowbreak{}error}}
\leansource{\texttt{correctedMeasure\_\allowbreak{}iterate\_\allowbreak{}weak\_\allowbreak{}error} \quad KineticLangevin/\allowbreak{}HamiltonianIteration.lean, lines 14–42}
\begin{leancode}
\leanline{0}{theorem\ correctedMeasure\_iterate\_weak\_error\ (π\ :\ Measure\ H)\ [IsFiniteMeasure\ π]}
\leanline{4}{(V\ R\ :\ H\ →\ H)\ (hV\ :\ Measurable\ V)\ (hR\ :\ Measurable\ R)}
\leanline{4}{(a\ :\ H\ →\ ℝ)\ (ha\ :\ Measurable\ a)\ (ha0\ :\ ∀\ z,\ 0\ ≤\ a\ z)\ (ha1\ :\ ∀\ z,\ a\ z\ ≤\ 1)}
\leanline{4}{(hstationary\ :\ (π.withDensity\ (fun\ z\ ↦\ ENNReal.ofReal\ (a\ z))).map\ V\ +}
\leanline{6}{(π.withDensity\ (fun\ z\ ↦\ ENNReal.ofReal\ (1-a\ z))).map\ R\ =\ π)}
\leanline{4}{(F\ :\ H\ →\ ℂ)\ (hF\ :\ Measurable\ F)\ (C\ :\ ℝ)\ (hC\ :\ 0\ ≤\ C)}
\leanline{4}{(hbound\ :\ ∀\ z,\ ‖F\ z‖\ ≤\ C)\ (n\ :\ ℕ)\ :}
\leanline{4}{‖(∫\ z,\ F\ (V\textasciicircum{}[n]\ z)\ ∂π)-(∫\ z,\ F\ z\ ∂π)‖\ ≤}
\leanline{6}{(n\ :\ ℝ)*(2*C)*(∫\ z,\ 1-a\ z\ ∂π)}
\end{leancode}
\leanlabel{How the Lean proof follows the argument}
\texttt{corrected\allowbreak{}Measure\_\allowbreak{}iterate\_\allowbreak{}weak\_\allowbreak{}error} proceeds by \texttt{induction n with}. The base cases establish the initial claims and the inductive branches use the induction hypotheses. Named proof applications include \texttt{hs}; the adjacent proof shows their hypotheses and order of use.
\end{leanbox}

\section{\texorpdfstring{The limiting pushforward through characteristic tests}{The limiting pushforward through characteristic tests}}


Suppose measurable maps \(Y_n\) converge pointwise to a measurable map \(Y\), and the bounded-test error is eventually at most a real sequence \(\delta_n\to0\). For every continuous real linear functional \(L\), use \(F(z)=\exp(iLz)\), whose norm is one. Dominated convergence under the finite measure gives \(\int F(Y_nz)\,d\pi\to\int F(Yz)\,d\pi\). The error bound gives convergence of the same sequence to \(\int F\,d\pi\). Uniqueness of limits equates the dual characteristic functions of \(Y_\#\pi\) and \(\pi\); the uniqueness theorem for finite measures gives \(Y_\#\pi=\pi\). The argument uses dual characteristic functions, so the result applies to Banach phase space with its usual product norm and does not silently require that norm to arise from an inner product.



Characteristic functions determine probability laws in finite-dimensional Euclidean space. If the numerical maps converge to the Hamiltonian flow and the accumulated test error vanishes, bounded convergence identifies the limiting pushforward with the target. This gives invariance of the flow.

\begin{leanbox}{Lean: \texttt{measure\_\allowbreak{}map\_\allowbreak{}eq\_\allowbreak{}of\_\allowbreak{}bounded\_\allowbreak{}test\_\allowbreak{}error\_\allowbreak{}tendsto}}
\leansource{\texttt{measure\_\allowbreak{}map\_\allowbreak{}eq\_\allowbreak{}of\_\allowbreak{}bounded\_\allowbreak{}test\_\allowbreak{}error\_\allowbreak{}tendsto} \quad KineticLangevin/\allowbreak{}HamiltonianIteration.lean, lines 48–78}
\begin{leancode}
\leanline{0}{theorem\ measure\_map\_eq\_of\_bounded\_test\_error\_tendsto}
\leanline{4}{(π\ :\ Measure\ H)\ [IsFiniteMeasure\ π]\ (Y\ :\ H\ →\ H)\ (hY\ :\ Measurable\ Y)}
\leanline{4}{(Yn\ :\ ℕ\ →\ H\ →\ H)\ (hYn\ :\ ∀\ n,\ Measurable\ (Yn\ n))}
\leanline{4}{(hconv\ :\ ∀\ z,\ Tendsto\ (fun\ n\ ↦\ Yn\ n\ z)\ atTop\ ({\leanfallback 𝓝}\ (Y\ z)))}
\leanline{4}{(δ\ :\ ℕ\ →\ ℝ)\ (hδ\ :\ Tendsto\ δ\ atTop\ ({\leanfallback 𝓝}\ 0))}
\leanline{4}{(hweak\ :\ ∀\ (F\ :\ H\ →\ ℂ),\ Measurable\ F\ →\ (∀\ z,\ ‖F\ z‖\ ≤\ 1)\ →}
\leanline{6}{∀ᶠ\ n\ in\ atTop,\ ‖(∫\ z,\ F\ (Yn\ n\ z)\ ∂π)-(∫\ z,\ F\ z\ ∂π)‖\ ≤\ δ\ n)\ :}
\leanline{4}{π.map\ Y\ =\ π}
\end{leancode}
\leanlabel{How the Lean proof follows the argument}
\texttt{measure\_\allowbreak{}map\_\allowbreak{}eq\_\allowbreak{}of\_\allowbreak{}bounded\_\allowbreak{}test\_\allowbreak{}error\_\allowbreak{}tendsto} gives the formal counterpart of this step. The main intermediate facts are \texttt{F}, \texttt{h\allowbreak{}F}, \texttt{h\allowbreak{}Fb}, \texttt{hactual}, \texttt{herr}, \texttt{htarget} in that order. Named proof applications include \texttt{tendsto\_\allowbreak{}integral\_\allowbreak{}of\_\allowbreak{}dominated\_\allowbreak{}convergence}, \texttt{ae\_\allowbreak{}of\_\allowbreak{}all}, \texttt{he}; the adjacent proof shows their hypotheses and order of use.
\end{leanbox}

\end{document}
